%!TEX root = ../template.tex
%%%%%%%%%%%%%%%%%%%%%%%%%%%%%%%%%%%%%%%%%%%%%%%%%%%%%%%%%%%%%%%%%%%%
%% abstract-pt.tex
%% NOVA thesis document file
%%
%% Abstract in Portuguese
%%%%%%%%%%%%%%%%%%%%%%%%%%%%%%%%%%%%%%%%%%%%%%%%%%%%%%%%%%%%%%%%%%%%

\typeout{NT FILE abstract-pt.tex}

Os Tipos de Dados Replicados sem Conflitos constituem uma componente fundamental em sistemas distribuídos modernos e aplicações de utilizador, permitindo que os dados sejam atualizados de forma independente de uma entidade centralizada e concorrente em múltiplas réplicas, garantindo uma consistência eventual forte. No entanto, o desenvolvimento destes tipos de dados apresenta desafios significativos, particularmente no que diz respeito à eficiência de memória, modularidade e à necessidade de uma verificação formal rigorosa para garantir que estes sejam seguros e corretos por construção.

Embora ferramentas especializadas como o VeriFx tenham surgido para verificar CRDTs, estas apresentam limitações críticas que dificultam o desenvolvimento de estruturas de dados complexas. Especificamente, o VeriFx demonstra dificuldades com algoritmos recursivos, falha na validação do alinhamento semântico entre as políticas de resolução de conflitos e a intenção do programador, e apresenta instabilidade na verificação de objetos complexos, tais como os que utilizam vetores de versão. Estas limitações forçam frequentemente os programadores a recorrer a abstrações manuais que podem introduzir erros ou levar a interrupções por tempo limite no processo de verificação.

Para abordar estas questões, esta tese propõe uma framework integrada para o desenvolvimento modular e formal de CRDTs, combinando os benefícios das ferramentas VeriFx e Why3.

\todo[inline]{Apresentar os resultados (implicações e consequências) da solução.}

% Abstract keywords
\keywords{
  CRDTs \and
  Verificação Formal \and
  Why3 \and
  VeriFx
}
